\chapter{Cronograma de atividades}
\label{cap:cronograma}

As atividades previstas para o TCC2 serão desenvolvidas no primeiro
semestre de 2027. O cronograma contempla a implementação dos protótipos,
a preparação dos cenários conversacionais, a execução do experimento,
a análise dos resultados e a conclusão da monografia.

A redação e as revisões com o orientador ocorrerão ao longo do semestre.
Os períodos previstos para a defesa e a entrega final serão ajustados
conforme o calendário oficial da unidade curricular.

\begin{table}[htb]
    \centering
    \caption{Cronograma previsto para o TCC2 no primeiro semestre de 2027}
    \label{tab:cronograma-tcc2}
    \small
    \begin{tabular}{p{0.52\textwidth}cccccc}
        \hline
        \textbf{Atividade}
        & \textbf{Fev.} & \textbf{Mar.} & \textbf{Abr.}
        & \textbf{Maio} & \textbf{Jun.} & \textbf{Jul.} \\
        \hline
        Revisão da base de conhecimento e dos prompts
        & X & X & & & & \\

        Implementação dos protótipos e dos registros de execução
        & X & X & & & & \\

        Elaboração e revisão dos cenários conversacionais
        & X & X & & & & \\

        Integração das métricas e execução do teste-piloto
        & & X & X & & & \\

        Congelamento das configurações e execução do experimento
        & & & X & & & \\

        Consolidação e análise comparativa dos resultados
        & & & X & X & & \\

        Redação da monografia e revisões com o orientador
        & X & X & X & X & X & \\

        Preparação da apresentação e defesa pública
        & & & & & X & \\

        Correções solicitadas pela banca e entrega final
        & & & & & X & X \\
        \hline
    \end{tabular}
    \fonteproprioautor
\end{table}

A execução do experimento será iniciada após a conclusão do teste-piloto
e o congelamento dos prompts, dos cenários e das configurações.
A análise utilizará os registros obtidos nessa execução. Após a defesa,
serão realizadas as correções solicitadas pela banca e a entrega da
versão final.